
% 09. Redistribution of Power and Dissipation of Momentum: High-Resolution Social Smoothing.tex 

\section{권력의 재분배와 모멘텀 해소: 고해상도 사회적 평활화}
우리는 오랫동안 인류의 역사를 억압에서 해방으로 나아가는 '자유의 확대 과정'으로 이해해 왔다. 한 명의 왕이 모든 것을 지배하던 야만의 시대에서 깨어난 시민들이 참여하는 시대로 발전하였고, 참혹했던 노예제는 폐지되었으며, 여성은 참정권을 획득하여 민주주의와 공정성이라는 숭고한 가치를 이룩했다고 믿는다. 우리는 이러한 변화를 도덕적 진보이자 정의의 승리로 맹신하며 박수를 보낸다.

그러나 디스턴스(DISTANCE)의 서늘한 렌즈를 들이대면, 역사는 전혀 다른 웅장한 물리학적 골조를 드러낸다. 인류의 역사는 정의와 부정의의 싸움도, 선과 악의 대립도 아니다. 오히려 그것은 사회라는 거대한 에너지 구조 내부에서 형성된 모멘텀들이 맹렬하게 끓어오르다 폭발하며 해소되는 역학적 과정이며, 더 거시적으로는 우주 전체가 완벽한 정지를 향해 추락하고 있는 '거대한 평활화(Equalization)' 궤적의 한 층위일 뿐이다.

초기 인류 사회의 기원을 살펴보자. 인류는 국가나 법, 의료 체계가 없는 척박한 자연 한가운데 던져져 맹수의 위협과 식량 부족, 질병과 싸워야 했다. 그 시대의 가장 숭고한 가치는 고결한 자유나 평등이 아니라 오직 '생존' 그 자체였다. 생존하기 위해서는 집단이 하나의 유기체처럼 빠르게 움직여야 했고, 전쟁이나 위기 상황에서 의사결정은 극도로 신속해야만 했다.

이러한 가혹한 환경 속에서 권력이 특정 위치에 집중되는 것은 탐욕의 결과가 아니라, 생존을 위해 자연스럽게 빚어진 거대한 '사회적 모멘텀'의 형성이었다. DISTANCE의 관점에서 왕이나 귀족, 부족장은 하늘이 내린 특별한 존재가 아니다. 그들은 단지 사회 전체가 품고 있는 거대한 모멘텀을 특정 방향으로 정렬하기 위해 억지로 빚어낸 구조적 '결절점(Node)'에 불과하다. 권력을 집중시키는 것은 집단을 빠르게 움직이고 질서를 유지하는 데 극도로 효율적이었으나, 그 대가로 지배자와 피지배자 사이에는 아찔할 만큼 막대한 전위차, 즉 '거대한 에너지 갭(Energy Gap)'이 형성되고 만다.

이 거대한 에너지 갭이 가장 끔찍하고 극단적인 형태로 고착된 구조가 바로 '노예제'이다. 한쪽은 모든 권리를 독점하고 다른 한쪽은 아무런 권리도 갖지 못하는 이 구조는, 현대인의 도덕적 잣대로는 부당하기 짝이 없으나 수천 년간 문명을 지탱해 왔다. 그 이유는 이 극단적인 에너지 갭이 뿜어내는 거대한 사회적 모멘텀이 엄청난 생산력을 강제로 창출해 냈기 때문이다. 그러나 물리학의 법칙은 예외가 없다. 에너지 갭이 극단적으로 응축되어 거대한 모멘텀이 존재하는 한, 그것을 무너뜨리려는 '해소'의 작용은 필연적으로 발생한다. 수많은 노예 반란과 피비린내 나는 저항은 우연한 비극이나 사회의 실패가 아니다. 그것은 억눌린 모멘텀이 존재한다는 증거이자, 갭을 허물어뜨리려는 자연스러운 우주적 폭발이다.

이 지점에서 우리가 찬미해 마지않는 '시민혁명'의 낭만 역시 철저히 전복된다. 프랑스 혁명을 떠올려 보라. 왕과 귀족이 막대한 권력을 독점하고 시민은 배제된 구조 속에서, 수백 년간 지독하게 축적된 거대한 사회적 에너지 갭은 마침내 임계점을 넘어 폭발했다. 우리는 이것을 자유를 향한 숭고한 투쟁으로 포장하지만, DISTANCE의 관점에서 혁명은 갑자기 하늘에서 뚝 떨어진 정의의 실현이 아니다. 그것은 단지 오랫동안 억눌려 있던 사회적 포텐셜 에너지(모멘텀)가 자신의 갭을 박살 내고 새로운 구조를 향해 맹렬하게 쏟아져 내린 '모멘텀 해소의 열역학적 폭포수'에 지나지 않는다.

왕정에서 민주주의로의 발전도 완벽히 동일한 역학을 따른다. 민주주의는 단순히 가장 윤리적이고 정의로운 정치 제도가 아니다. 그것은 거대한 권력(모멘텀)을 수많은 개인에게 잘게 쪼개어 분산시키는 극도로 정교한 '모멘텀 분산 기술'이다. 한 명의 왕에게 집중되어 있던 거대한 에너지 갭을 의회로, 그리고 수천만 명의 시민에게로 분산시킴으로써, 사회는 거대한 덩어리가 부딪히는 '낮은 해상도의 사회'에서 수많은 작은 모멘텀들이 잘게 쪼개져 맞물려 돌아가는 '높은 해상도의 균질한 구조'로 진화해 간 것이다.

이 과정에서 매우 소름 돋는 역설이 발견된다. 흔히 민주주의와 평등이 사람들을 더 가깝게 하나로 묶었다고 생각하지만, 물리학적 실상은 정반대다. 왕과 노예가 존재하던 시대에는 둘 사이의 아찔한 에너지 갭 때문에 서로를 향한 맹렬한 모멘텀이 당겨지며 극도로 '짧은 디스턴스'가 유지되었다. 그러나 혁명과 민주주의가 이 거대한 갭을 무너뜨리고 모멘텀을 평탄하게 분산시켜 버리자, 사회 전체는 거대한 낙차가 소거되어 서로가 서로를 향해 맹렬히 부딪힐 필요가 없는 '더욱 긴 디스턴스'의 구조로 흩어지게 된 것이다.

결국 우리가 목격한 여성 참정권의 획득이나 공정성의 확대는 모두 단순한 선악의 승리가 아니다. 육체적 힘이 지배하던 생존 사회의 낡은 구조가 효력을 잃으면서, 남녀 사이에 존재하던 에너지 갭이 폭발하고 새로운 모멘텀 재배치가 일어난 것일 뿐이다. 공정성이란 특정 집단이 독점하던 거대한 모멘텀을 작은 모멘텀들로 쪼개어 사회의 전위차를 줄여나가는 기계적인 깎아냄의 과정이다.

인류의 역사는 왕이 사라지고 시민이 등장한 얄팍한 서사가 아니다. 그것은 사회라는 거대한 에너지 구조가 스스로 품고 있던 거대한 권력 모멘텀을 맹렬히 해소하며, 더 높은 해상도를 가진 완벽한 평활화(Equalization)를 향해 추락해 온 우주적 진화의 연장선이다. 생명이 물질 세계의 모멘텀을 깎아내는 로컬 평활 엔진이라면, 인류의 사회와 제도는 그 생명 위에 구축된 '가장 웅장하고 치밀한 사회적 평활화 엔진'인 것이다. 우리가 피와 땀으로 기록해 온 역사란, 결국 우주가 자신의 남은 에너지 갭을 모조리 깎아내기 위해 맹렬하게 작동시키고 있는 거대한 엔진의 시끄러운 파열음들에 불과하다.
